\documentclass[aspectratio=169,11pt]{beamer}
\usepackage[indonesian]{babel}
\usepackage{amsmath,amssymb,booktabs,listings,tikz,xcolor,hyperref}
\usetheme{metropolis}
\usefonttheme{professionalfonts}
\metroset{progressbar=frametitle,block=fill}
\usebeamercolor{normal text}\colorlet{codeink}{normal text.fg}
\usebeamercolor{alerted text}\colorlet{codeaccent}{alerted text.fg!60!codeink}
\lstset{language={[LaTeX]TeX},basicstyle=\ttfamily\scriptsize\color{codeink},
  texcsstyle=*\bfseries\color{codeaccent},columns=fullflexible,keepspaces=true,
  showstringspaces=false,backgroundcolor=\color{codeink!6},frame=l,
  rulecolor=\color{codeaccent},xleftmargin=5pt,aboveskip=4pt,belowskip=4pt,
  breaklines=true}
\lstdefinestyle{terminal}{language=bash,basicstyle=\ttfamily\scriptsize\color{codeink},
  keywordstyle=\bfseries\color{codeaccent},breaklines=true}
\lstdefinestyle{plain}{language={},basicstyle=\ttfamily\scriptsize\color{codeink},
  columns=fullflexible,keepspaces=true,breaklines=true}
\newcommand{\tugas}[1]{\par\medskip\alert{\textbf{Coba sendiri}}\quad #1}
\title{Membuat kelas dan paket \LaTeX}
\subtitle{Hari 4 --- Arsitektur, implementasi, dan distribusi}
\author{Pelatihan Universitas Terbuka dan Departemen Matematika FMIPA UGM}
\date{Materi lanjutan}

\begin{document}
\begin{frame}[plain]
\usebeamerfont{title}\inserttitle\par
\usebeamerfont{subtitle}\insertsubtitle\par\medskip
\usebeamertemplate{title separator}\par\medskip
\usebeamerfont{author}\insertauthor\par\smallskip
\usebeamerfont{date}\insertdate
\end{frame}

\begin{frame}{Tujuan belajar Hari 4}
\begin{enumerate}
  \item Membedakan tanggung jawab dokumen, kelas \texttt{.cls}, dan paket \texttt{.sty}.
  \item Menulis paket kecil yang memiliki opsi, perintah, dan environment yang dapat digunakan ulang.
  \item Membuat kelas dokumen yang mewarisi kelas standar lalu menambahkan aturan institusi.
  \item Menguji, mendokumentasikan, dan mengemas hasil agar mudah dipakai orang lain.
\end{enumerate}
\end{frame}

\begin{frame}{Peta belajar Hari 4}
\tableofcontents
\end{frame}

\section{Sesi 1. Arsitektur kelas dan paket}

\begin{frame}{Sesi 1 --- Arsitektur kelas dan paket}
\large 08.30--10.00 \quad 90 menit
\medskip
\begin{itemize}
  \item 20 menit: hubungan dokumen, kelas, paket, dan toolchain
  \item 25 menit: anatomi berkas \texttt{.sty} dan \texttt{.cls}
  \item 25 menit: pewarisan kelas dan pemrosesan opsi
  \item 20 menit: latihan membaca template dan memetakan dependensi
\end{itemize}
\end{frame}

\begin{frame}{Empat lapisan yang perlu dipisahkan}
\begin{columns}[T,onlytextwidth]
\begin{column}{.23\textwidth}\centering
  \textbf{Dokumen}\par\smallskip
  Isi, struktur, dan metadata.
\end{column}
\begin{column}{.23\textwidth}\centering
  \textbf{Kelas}\par\smallskip
  Aturan dasar jenis dokumen.
\end{column}
\begin{column}{.23\textwidth}\centering
  \textbf{Paket}\par\smallskip
  Fitur tambahan yang dapat dipilih.
\end{column}
\begin{column}{.23\textwidth}\centering
  \textbf{Mesin}\par\smallskip
  LaTeX, XeLaTeX, atau LuaLaTeX.
\end{column}
\end{columns}
\bigskip
\begin{block}{Prinsip desain}
Simpan keputusan yang berulang pada lapisan yang paling tepat. Isi jangan
mengetahui detail warna institusi, dan kelas jangan memaksa fitur yang hanya
dibutuhkan satu dokumen.
\end{block}
\end{frame}

\begin{frame}[fragile]{Memilih kelas atau paket}
\begin{lstlisting}
\documentclass[12pt,a4paper]{article} % memilih jenis dokumen
\usepackage{utstyle}                   % menambah fitur atau gaya
\begin{document}
Isi dokumen tetap berada di berkas utama.
\end{document}
\end{lstlisting}
\begin{itemize}
  \item Gunakan kelas untuk aturan global: ukuran halaman, struktur, judul, dan heading.
  \item Gunakan paket untuk fitur yang dapat dipakai lintas kelas.
  \item Gunakan prambule biasa untuk keputusan lokal pada satu dokumen.
\end{itemize}
\tugas{Temukan tiga baris pada dokumen peserta yang seharusnya dipindahkan ke paket.}
\end{frame}

\begin{frame}[fragile]{Anatomi minimal paket \texttt{.sty}}
\begin{lstlisting}
\NeedsTeXFormat{LaTeX2e}
\ProvidesPackage{utstyle}[2026/11/19 v0.1 Gaya dokumen UT]
\RequirePackage{xcolor}
\newcommand{\institusi}{Universitas Terbuka}
\endinput
\end{lstlisting}
\begin{description}
  \item[\texttt{\textbackslash NeedsTeXFormat}] menyatakan format minimum yang diharapkan.
  \item[\texttt{\textbackslash ProvidesPackage}] memberi nama, tanggal, dan versi.
  \item[\texttt{\textbackslash RequirePackage}] memuat dependensi paket.
  \item[\texttt{\textbackslash endinput}] menandai akhir kode yang diproses.
\end{description}
\end{frame}

\begin{frame}[fragile]{Anatomi minimal kelas \texttt{.cls}}
\begin{lstlisting}
\NeedsTeXFormat{LaTeX2e}
\ProvidesClass{utreport}[2026/11/19 v0.1 Kelas laporan UT]
\LoadClass[12pt,a4paper]{article}
\RequirePackage{geometry}
\geometry{margin=3cm}
\endinput
\end{lstlisting}
Kelas biasanya memuat kelas lain dengan \lstinline|\LoadClass|, kemudian
menambahkan aturan yang menjadi identitas seluruh dokumen. Hindari menyalin
isi \texttt{article.cls} atau \texttt{book.cls} ke berkas sendiri.
\end{frame}

\begin{frame}[fragile]{Waktu pemrosesan perintah}
\begin{lstlisting}
\documentclass{utreport}  % kelas dibaca lebih dahulu
\usepackage{utstyle}      % paket dibaca sesudah kelas
\begin{document}          % isi dokumen dimulai
\end{lstlisting}
\begin{itemize}
  \item Kelas dapat menentukan opsi dan memuat dependensi dasar.
  \item Paket dapat dipanggil setelah kelas, sehingga tidak boleh mengasumsikan isi dokumen sudah dimulai.
  \item Perintah yang harus dijalankan setelah semua paket tersedia dapat ditempatkan pada hook.
\end{itemize}
\end{frame}

\begin{frame}[fragile]{Pewarisan kelas dengan \texttt{\textbackslash LoadClass}}
\begin{lstlisting}
\DeclareOption{draft}{\PassOptionsToClass{draft}{article}}
\DeclareOption*{\PassOptionsToClass{\CurrentOption}{article}}
\ProcessOptions\relax
\LoadClass[12pt,a4paper]{article}
\end{lstlisting}
\begin{itemize}
  \item Teruskan opsi yang memang dimengerti kelas induk.
  \item Jangan menerima opsi diam-diam jika opsi itu mengubah hasil secara penting.
  \item Dokumentasikan opsi publik dan perilaku defaultnya.
\end{itemize}
\end{frame}

\begin{frame}[fragile]{Opsi paket dengan \texttt{\textbackslash DeclareOption}}
\begin{lstlisting}
\newif\ifut@compact
\DeclareOption{compact}{\ut@compacttrue}
\DeclareOption{spacious}{\ut@compactfalse}
\DeclareOption*{\PackageWarning{utstyle}{Opsi tidak dikenal: \CurrentOption}}
\ProcessOptions\relax
\ifut@compact
  \setlength{\parskip}{3pt}
\else
  \setlength{\parskip}{8pt}
\fi
\end{lstlisting}
Opsi adalah antarmuka publik. Nama opsi perlu stabil, sedangkan cara
implementasinya boleh berubah di versi berikutnya.
\end{frame}

\begin{frame}{Kontrak antarmuka sebuah paket}
\begin{columns}[T,onlytextwidth]
\begin{column}{.47\textwidth}
\textbf{Publik}\par\smallskip
Perintah, environment, opsi, dan nama warna yang sengaja didokumentasikan.
\end{column}
\begin{column}{.47\textwidth}
\textbf{Internal}\par\smallskip
Counter, flag, helper, dan nama teknis yang dapat berubah tanpa memecahkan dokumen.
\end{column}
\end{columns}
\bigskip
Gunakan awalan nama yang konsisten, misalnya \texttt{\textbackslash ut...}, untuk mengurangi
risiko benturan dengan paket lain. Jangan membuat ulang perintah umum tanpa
memeriksa apakah nama tersebut sudah dipakai.
\end{frame}

\begin{frame}{Latihan pemetaan dependensi}
Ambil satu template yang pernah digunakan, lalu buat tabel kecil:
\begin{enumerate}
  \item nama berkas utama, kelas, dan semua paket;
  \item perintah atau environment yang disediakan tiap paket;
  \item bagian yang hanya mengatur gaya institusi;
  \item opsi dan mesin kompilasi yang diperlukan.
\end{enumerate}
\tugas{Tandai satu aturan yang sebaiknya menjadi kelas dan satu aturan yang sebaiknya tetap menjadi paket.}
\end{frame}

\section{Sesi 2. Menulis paket reusable}

\begin{frame}{Sesi 2 --- Menulis paket reusable}
\large 10.15--11.45 \quad 90 menit
\medskip
\begin{itemize}
  \item 15 menit: memilih antarmuka dan nama perintah
  \item 25 menit: membangun paket \texttt{utstyle.sty} secara bertahap
  \item 25 menit: commands, environments, warna, dan konfigurasi
  \item 25 menit: latihan menguji paket pada dua dokumen berbeda
\end{itemize}
\end{frame}

\begin{frame}[fragile]{Mulai dari antarmuka, bukan dari dekorasi}
Tuliskan kebutuhan pengguna dalam bentuk yang dapat diuji.
\begin{lstlisting}
\usepackage[blue]{utstyle}
\namauniversitas{Universitas Terbuka}
\begin{tujuan}
Peserta mampu membuat paket yang dapat digunakan ulang.
\end{tujuan}
\end{lstlisting}
Kemudian tentukan: apakah perintah menerima argumen, kapan dipanggil, apa
hasil defaultnya, dan pesan apa yang muncul jika pengguna salah memakainya.
\end{frame}

\begin{frame}[fragile]{Paket gaya kecil: metadata dan dependensi}
\begin{lstlisting}
\NeedsTeXFormat{LaTeX2e}
\ProvidesPackage{utstyle}[2026/11/19 v0.1 Gaya pelatihan]
\RequirePackage{xcolor}
\RequirePackage{etoolbox}
\definecolor{utblue}{HTML}{005B96}
\newcommand{\namauniversitas}[1]{\def\ut@university{#1}}
\namauniversitas{Universitas Terbuka}
\end{lstlisting}
Nama internal memakai awalan khusus. Untuk memakai karakter \texttt{@} dalam
definisi yang benar-benar dieksekusi, gunakan \lstinline|\makeatletter| dan
\lstinline|\makeatother| bila diperlukan.
\end{frame}

\begin{frame}[fragile]{Perintah dengan argumen wajib dan opsional}
\begin{lstlisting}
\newcommand{\hasil}[1]{\textbf{Hasil:} #1}
\newcommand{\catatan}[2][Catatan]{\textit{#1:} #2}
\hasil{Periode bertambah ketika tali diperpanjang.}
\catatan{Periksa satuan sebelum menghitung.}
\catatan[Perhatian]{Jangan campur cm dan m.}
\end{lstlisting}
Pilih perintah yang menyatakan makna, bukan sekadar tampilan. Jika definisi
mulai memiliki banyak cabang, pertimbangkan environment atau key--value.
\end{frame}

\begin{frame}[fragile]{Environment untuk pola isi berulang}
\begin{lstlisting}
\newenvironment{tujuan}
  {\begin{block}{Tujuan sesi}}
  {\end{block}}
\newenvironment{peringatan}
  {\begin{alertblock}{Perhatian}}
  {\end{alertblock}}
\end{lstlisting}
Environment mengelompokkan awal dan akhir suatu pola. Pengguna cukup menulis
isi semantis, sedangkan pilihan tampilan dipusatkan pada paket.
\end{frame}

\begin{frame}[fragile]{Pengaturan warna dan gaya}
\begin{lstlisting}
\definecolor{utblue}{HTML}{005B96}
\colorlet{utaccent}{utblue!80!black}
\newcommand{\judulwarna}[1]{\textcolor{utaccent}{\textbf{#1}}}
\setlength{\parindent}{0pt}
\setlength{\parskip}{6pt plus 2pt minus 1pt}
\end{lstlisting}
Simpan angka, warna, dan nama gaya di satu tempat. Hindari memaksa font atau
ukuran yang bertentangan dengan kelas dan kebutuhan pembaca.
\end{frame}

\begin{frame}[fragile]{Hook untuk konfigurasi setelah paket dimuat}
\begin{lstlisting}
\AtBeginDocument{%
  \providecommand{\judulpelatihan}{Pelatihan LaTeX}%
}
\AtBeginEnvironment{figure}{\centering}
\end{lstlisting}
Gunakan hook dengan hati-hati. Dependensi antar-paket yang tidak jelas akan
menyulitkan diagnosis. Jika urutan penting, dokumentasikan alasannya.
\end{frame}

\begin{frame}[fragile]{Validasi input sederhana}
\begin{lstlisting}
\newcommand{\tahun}[1]{%
  \ifnum#1<2000
    \PackageError{utstyle}{Tahun terlalu kecil}{Gunakan tahun 2000 atau lebih.}%
  \else
    \def\ut@year{#1}%
  \fi
}
\end{lstlisting}
Pesan error yang baik menyebutkan apa yang salah dan bagaimana memperbaikinya.
Validasi secukupnya: jangan membuat paket menjadi lebih rumit daripada
masalah yang hendak diselesaikan.
\end{frame}

\begin{frame}[fragile]{Paket yang sudah dirakit}
\begin{lstlisting}
% utstyle.sty
\NeedsTeXFormat{LaTeX2e}
\ProvidesPackage{utstyle}[2026/11/19 v0.1 Gaya pelatihan]
\RequirePackage{xcolor}
\definecolor{utblue}{HTML}{005B96}
\newcommand{\hasil}[1]{\textcolor{utblue}{\textbf{Hasil:} #1}}
\newenvironment{tujuan}{\begin{quote}\textbf{Tujuan.}\ }{\end{quote}}
\endinput
\end{lstlisting}
\begin{lstlisting}
% main.tex
\documentclass{article}
\usepackage{utstyle}
\begin{document}
\begin{tujuan}Membuat paket yang dapat digunakan ulang.\end{tujuan}
\hasil{Paket berhasil dimuat.}
\end{document}
\end{lstlisting}
\end{frame}

\begin{frame}{Pengujian pada dokumen minimal}
Buat dokumen uji sekecil mungkin agar sumber masalah cepat terlihat:
\begin{enumerate}
  \item satu kelas standar;
  \item satu paket yang sedang dikembangkan;
  \item satu contoh untuk setiap perintah dan environment;
  \item satu contoh input yang salah untuk memeriksa pesan error.
\end{enumerate}
\tugas{Kompilasi dengan kelas \texttt{article} dan \texttt{book}. Catat perilaku yang memang berbeda.}
\end{frame}

\begin{frame}[fragile]{Uji cepat dan pemeriksaan log}
\begin{lstlisting}[style=terminal]
latexmk -pdf -interaction=nonstopmode test-utstyle.tex
rg "Undefined|Error|Warning|Overfull" test-utstyle.log
\end{lstlisting}
Uji PDF, bukan hanya keberhasilan kompilasi. Periksa judul, margin, warna,
nomor, rujukan, dan hasil saat paket dimuat bersama paket lain.
\end{frame}

\begin{frame}{Rilis kecil yang dapat dipercaya}
Sebelum membagikan paket:
\begin{itemize}
  \item naikkan versi jika antarmuka berubah;
  \item tulis daftar dependensi dan versi minimum;
  \item sertakan contoh minimal dan cara kompilasi;
  \item bedakan error yang harus menghentikan kompilasi dari warning informatif;
  \item catat perubahan pada berkas changelog.
\end{itemize}
\end{frame}

\section{Sesi 3. Menulis kelas dan mendistribusikannya}

\begin{frame}{Sesi 3 --- Menulis kelas dan mendistribusikannya}
\large 13.00--14.30 \quad 90 menit
\medskip
\begin{itemize}
  \item 20 menit: merancang kelas di atas \texttt{article} atau \texttt{report}
  \item 25 menit: judul, heading, metadata, dan opsi kelas
  \item 25 menit: proyek mini \texttt{utreport.cls}
  \item 20 menit: pengujian, dokumentasi, dan paket distribusi
\end{itemize}
\end{frame}

\begin{frame}[fragile]{Kapan membuat kelas baru?}
\begin{itemize}
  \item Buat paket jika fitur dapat dipakai pada beberapa jenis dokumen.
  \item Buat kelas jika struktur dan aturan global memang berbeda.
  \item Perluas kelas standar jika sebagian besar perilakunya masih sesuai.
\end{itemize}
\begin{lstlisting}
\ProvidesClass{utreport}[2026/11/19 v0.1 Kelas laporan]
\LoadClass[12pt,a4paper]{report}
\RequirePackage{geometry}
\RequirePackage{fancyhdr}
\geometry{margin=3cm}
\end{lstlisting}
Mulai dari pewarisan. Menulis kelas dari nol membuat banyak keputusan
internal yang tidak perlu dan meningkatkan biaya pemeliharaan.
\end{frame}

\begin{frame}[fragile]{Opsi kelas yang diteruskan}
\begin{lstlisting}
\DeclareOption{twoside}{\PassOptionsToClass{twoside}{report}}
\DeclareOption{draft}{\PassOptionsToClass{draft}{report}}
\DeclareOption{final}{\PassOptionsToClass{final}{report}}
\DeclareOption*{\PassOptionsToClass{\CurrentOption}{report}}
\ProcessOptions\relax
\LoadClass[12pt,a4paper]{report}
\end{lstlisting}
Jangan memanggil \lstinline|\LoadClass| sebelum opsi diproses jika opsi itu
harus diteruskan ke kelas induk. Buat daftar opsi publik dalam dokumentasi.
\end{frame}

\begin{frame}[fragile]{Mengatur metadata dan halaman judul}
\begin{lstlisting}
\newcommand{\judulresmi}[1]{\def\ut@title{#1}}
\newcommand{\penulisresmi}[1]{\def\ut@author{#1}}
\renewcommand{\maketitle}{%
  \begin{titlepage}
    \centering
    {\Large\bfseries \ut@university\par\vfill}
    {\huge\bfseries \ut@title\par\medskip}
    {\large \ut@author\par\vfill}
  \end{titlepage}}
\end{lstlisting}
Sediakan default yang aman atau berikan error yang jelas jika metadata wajib
belum diisi. Jangan mengandalkan urutan internal yang tidak terdokumentasi.
\end{frame}

\begin{frame}[fragile]{Heading dan gaya halaman}
\begin{lstlisting}
\RequirePackage{titlesec}
\titleformat{\chapter}[display]
  {\normalfont\bfseries\Large}
  {\chaptertitlename\ \thechapter}{8pt}{\Huge}
\RequirePackage{fancyhdr}
\pagestyle{fancy}
\fancyhf{}
\fancyhead[LE,RO]{\thepage}
\fancyhead[LO]{\nouppercase{\rightmark}}
\fancyhead[RE]{\nouppercase{\leftmark}}
\end{lstlisting}
Pastikan perubahan berlaku pada halaman biasa, halaman pembuka bab, dan mode
\texttt{twoside}. Uji juga judul panjang dan halaman kosong.
\end{frame}

\begin{frame}[fragile]{Kelas mini \texttt{utreport.cls}}
\begin{lstlisting}
\NeedsTeXFormat{LaTeX2e}
\ProvidesClass{utreport}[2026/11/19 v0.1 Kelas laporan UT]
\DeclareOption*{\PassOptionsToClass{\CurrentOption}{report}}
\ProcessOptions\relax
\LoadClass[12pt,a4paper]{report}
\RequirePackage{geometry}
\RequirePackage{utstyle}
\geometry{margin=3cm}
\setlength{\parindent}{0pt}
\setlength{\parskip}{6pt plus 2pt minus 1pt}
\endinput
\end{lstlisting}
Kelas ini memusatkan keputusan yang berlaku untuk semua dokumen laporan.
Fitur semantis yang dapat dipakai lintas kelas tetap berada di \texttt{utstyle}.
\end{frame}

\begin{frame}[fragile]{Memakai kelas mini}
\begin{lstlisting}
\documentclass[oneside]{utreport}
\title{Laporan Eksperimen Ayunan}
\author{Nama Peserta}
\date{19 November 2026}
\begin{document}
\maketitle
\tableofcontents
\chapter{Pendahuluan}
Pertanyaan penelitian dan konteks eksperimen.
\end{document}
\end{lstlisting}
\begin{block}{Kriteria selesai}
Dokumen pengguna hanya memilih kelas dan menulis isi. Detail margin,
indentasi, warna, serta dependensi tidak diulang di setiap proyek.
\end{block}
\end{frame}

\begin{frame}[fragile]{Struktur proyek untuk distribusi}
\begin{columns}[T,onlytextwidth]
\begin{column}{.43\textwidth}
\begin{lstlisting}[style=plain]
utreport/
  utreport.cls
  utstyle.sty
  README.md
  CHANGELOG.md
  example.tex
  LICENSE
\end{lstlisting}
\end{column}
\begin{column}{.52\textwidth}
\begin{itemize}
  \item README menjelaskan pemasangan dan kompilasi.
  \item Contoh menunjukkan antarmuka publik, bukan detail internal.
  \item Changelog mencatat perubahan yang terlihat pengguna.
  \item Lisensi mengikuti sumber dan kebijakan institusi.
\end{itemize}
\end{column}
\end{columns}
\end{frame}

\begin{frame}[fragile]{Kompilasi bersih dan mesin kompilasi}
\begin{lstlisting}[style=terminal]
latexmk -C
latexmk -pdf example.tex
% atau, untuk font dan bahasa modern:
latexmk -xelatex example.tex
\end{lstlisting}
Uji dari folder proyek baru agar paket tidak kebetulan ditemukan dari folder
kerja pribadi. Jika kelas akan dipasang secara global, dokumentasikan lokasi
dan cara menghapus versi lama.
\end{frame}

\begin{frame}{Checklist API dan kualitas}
\begin{enumerate}
  \item Nama kelas dan paket memiliki versi serta tanggal.
  \item Opsi, perintah, environment, dan dependensi terdokumentasi.
  \item Nama internal tidak bocor sebagai API publik.
  \item Dokumen minimal berhasil dikompilasi dengan kelas dan mesin target.
  \item Pesan error membantu pengguna memperbaiki sumber.
  \item Contoh dan changelog diperbarui bersama kode.
\end{enumerate}
\end{frame}

\begin{frame}{Proyek akhir Hari 4}
Buat paket atau kelas kecil untuk kebutuhan nyata peserta:
\begin{enumerate}
  \item tulis satu paragraf kebutuhan pengguna;
  \item definisikan API minimal dengan paling banyak tiga perintah publik;
  \item buat satu berkas \texttt{.sty} atau \texttt{.cls} dan satu contoh;
  \item uji dokumen pada kompilasi bersih;
  \item minta rekan menggunakan contoh tanpa melihat implementasinya.
\end{enumerate}
\tugas{Serahkan sumber, PDF contoh, README, dan catatan batasan versi pertama.}
\end{frame}

\begin{frame}{Ringkasan Hari 4}
\begin{itemize}
  \item Kelas mengatur kerangka dokumen; paket menyediakan fitur yang dapat digunakan ulang.
  \item Antarmuka publik perlu kecil, jelas, berversi, dan diuji.
  \item Pewarisan kelas standar biasanya lebih aman daripada menulis dari nol.
  \item Distribusi yang baik mencakup contoh, dokumentasi, changelog, dan kompilasi bersih.
\end{itemize}
\bigskip
\centering\large Pertanyaan dan konsultasi
\end{frame}

\end{document}
